\section{Conclusions}\label{sec:conclusions}

On the bi-objective CDP and GDP, a MOEA/D framework with a pool of eight operators and normalised objectives is significantly better than untuned NSGA-II, GRASP and VNS baselines, at sizes up to 500 and under three quality indicators. It stays ahead of tuned NSGA-II, SMS-EMOA and GRASP and of a tabu search and a VNS in the style of Colmenar et al. when all methods are tuned with the same protocol and run for the same wall-clock time --- clearly on the GDP, and on the CDP up to the point where five times the time lets tuned GRASP, VNS and tabu search catch up. On the unconstrained MDP the order reverses: the tabu search, the VNS and a tuned GRASP are better than MOEA/D, which falls short at both ends of the front and does not improve with five times the time on the small instances; seeding the MOEA/D with the archive of the VNS for half of the time removes most of the gap (0.971 against 0.977 for the VNS with 100 to 150 elements) without costing anything on the CDP and GDP with n $\leq$ 150 \todo{n = 500 pending}.

Adaptive operator selection does not improve on uniform selection from the pool on instances up to about 200 elements nor on the CDP, whatever the selector (UCB1, Thompson sampling, the published DYTS, FRRMAB), reward, hyper-parameters, budget or context, and the scalar-reward selectors are significantly worse. The HV-reward selectors do win on large instances (twenty of twenty MDG-b instances with 2,000 elements), on tightly constrained ones (GDP with the 0.2 cost factor), with an advantage that grows as the budget shrinks, and, when the pool holds operators that choose the size of the set, on the CDP and GDP with 500 elements (up to +0.26); there the search is far from convergence at the standard budget. The pool needs MaxMinSwap, is robust to the removal of its three costly or redundant operators and tolerates three deliberately useless ones. A short MaxMin intensification brings the front's MaxMin end to within about 1\% of the best published CDP values on four of five instance families (0.99 on GKD-c; 0.85 on the hardest, MDG-b), and above the published GDP values. In the CDP and GDP the size of the set is a hidden dimension of the bi-objective problem, which the standard methods leave to the drift of the repair step (about twice the smallest feasible size on the CDP); two operators that choose the size raise the hypervolume of the MOEA/D several-fold, a cardinality bound makes the front finite, and on the CDP the ranking of the methods then depends on the bound. We release data and code to let others test selection policies and baselines against the uniform pool we propose.
